\documentclass[10pt,a4paper,sans]{moderncv}

% Estilo del CV (Classic, Banking, Casual...)
\moderncvstyle{classic} 
\moderncvcolor{blue} % Colores: blue, orange, green, red, purple, grey, black

\usepackage[utf8]{inputenc}
\usepackage[scale=0.75]{geometry}

% Datos Personales
\name{M. del Rocío}{Alarcón}
\address{San Rafael 250 - Las Heras}{Mendoza}
\phone[mobile]{261-4602671}
\email{alarcon.rocio.1h1@gmail.com}
\extrainfo{Nacimiento: 17 de septiembre del 2001}

% Foto (Asegúrate de subir el archivo como foto.jpg)
\photo[64pt][0.4pt]{foto.jpg}

\begin{document}
\makecvtitle

\section{Perfil Profesional}
\cvitem{}{Estudiante avanzado de Ingeniería Industrial (4to año) con sólida capacidad analítica y enfoque en la mejora de procesos. Cuento con un perfil bilingüe (Inglés y Portugués) y experiencia comprobable en atención al cliente y gestión administrativa. Me destaco por mi proactividad, rápida capacidad de aprendizaje y una fuerte orientación a la resolución de problemas en entornos dinámicos.}

\section{Experiencia Laboral}
\cventry{Oct 2025--Act.}{Administración y ventas}{Valentine Cocina}{}{}{
\begin{itemize}
    \item Gestión de ventas y atención personalizada al público.
    \item Tareas administrativas y control de stock.[cite: 1]
    \item Resolución de consultas y gestión de pedidos.
\end{itemize}}

\cventry{Jun 2022--Dic 2022}{Producción y ventas}{SuperMatize}{}{}{
\begin{itemize}
    \item Elaboración y manufactura de productos.[cite: 1]
    \item Atención al público y cierre de caja.
    \item Optimización de tareas manuales en el proceso productivo.
\end{itemize}}

\cventry{Eventual}{Staff de eventos}{Independiente}{}{}{
\begin{itemize}
    \item Adaptabilidad a ambientes de alta exigencia y trabajo bajo presión.
    \item Manejo de relaciones interpersonales y servicio al cliente.
\end{itemize}}

\section{Formación Académica}
\cventry{2020--En curso}{Ingeniería Industrial}{Universidad Nacional de Cuyo}{}{}{}
\cventry{2015--2019}{Bachiller en Cs. Sociales y Humanidades}{Escuela del Magisterio}{}{}{}

\section{Idiomas y Cursos}
\cvitem{Inglés}{Nivel B2 - Facultad de Filosofía y Letras, UNCuyo (2022).}
\cvitem{Portugués}{Nivel B2 - Facultad de Filosofía y Letras, UNCuyo (2020-2022).}
\cvitem{Excel}{Nivel Intermedio - Facultad de Ingeniería, UNCuyo (2020).}
\cvitem{Inglés}{Nivel B1 - Colegio de Lenguas Extranjeras, UNCuyo (2014-2018).}

\end{document}